% GENERATED FILE. Edit canonical lessons, then run npm run generate:books.
% canonical-source-hash: fnv1a64:50cc8a7fd8a4de31
% canonical-lessons: UR-C327-khazana, UR-C327-garna, UR-C327-nag, UR-C327-dasna, UR-C327-hira

\chapter{Buried Treasure}
\label{ch:ur-buried-treasure}
\hlchaptermodality{327}

\begin{chapteropening}
\textbf{By the end of this chapter:} \emph{I can tell a tale of hidden treasure.}

Complete the last lesson of chapter 327: I can tell a tale of hidden treasure.
\end{chapteropening}

\section[\texorpdfstring{\ur{خزانہ} (khazānā)}{خزانہ (khazānā)}]{\ur{خزانہ} (khazānā) — a treasure; a treasury}

\label{lesson:UR-C327-khazana}



\begin{quote}
Before the new one: say the Urdu for to sulk, then the Urdu for torment.
\end{quote}

\subsection*{You'll want to know: \ur{خزانہ}}
\textbf{\ur{خزانہ}} — \emph{khazānā} — \textquotedblleft{}a treasure; a treasury\textquotedblright{}.

\begin{quote}
\ur{زمین} \ur{کے} \ur{نیچے} \ur{ایک} \ur{خزانہ} \ur{تھا۔}
\end{quote}

\begin{quote}
\emph{zamīn ke nīche ek khazānā thā.}
\end{quote}

\begin{quote}
There was a treasure under the ground.
\end{quote}

Masculine.

\begin{cousinweb}
\textbf{\ur{خزانہ}} (\emph{khazānā}) is usually traced, through Persian, to Arabic \emph{khizāna}, \textquotedblleft{}a storehouse\textquotedblright{}. \textbf{\ur{خزانچی}}, a cashier, is the one who keeps it.
\end{cousinweb}

\begin{grammarlens}[title={What you've built}]
One more word for telling what happened.
\end{grammarlens}

\subsection*{Your turn}
\begin{itemize}
\raggedright
  \item \emph{Say it:} \emph{khazānā}
  \item \emph{Say it:} \emph{khazānā}, once more
  \item \emph{Say it:} \emph{zamīn ke nīche ek khazānā thā.} — There was a treasure under the ground.
  \item \emph{Recall:} say the Urdu for to sulk, then the Urdu for torment, then say \emph{khazānā} again
\end{itemize}

\begin{tcolorbox}[breakable,colback=teal!4,colframe=teal!35!black,title={Before you move on}]
What does \ur{خزانہ} mean? (\textquotedblleft{}A treasure; a treasury\textquotedblright{}.) Say it once more.
\end{tcolorbox}
\section[\texorpdfstring{\ur{گاڑنا} (gāṛnā)}{گاڑنا (gāṛnā)}]{\ur{گاڑنا} (gāṛnā) — to bury}

\label{lesson:UR-C327-garna}



\begin{quote}
Before the new one: say the Urdu for torment, then the Urdu for a treasure.
\end{quote}

\subsection*{You'll want to know: \ur{گاڑنا}}
\textbf{\ur{گاڑنا}} — \emph{gāṛnā} — \textquotedblleft{}to bury\textquotedblright{}.

\begin{quote}
\ur{چور} \ur{نے} \ur{خزانہ} \ur{زمین} \ur{میں} \ur{گاڑا۔}
\end{quote}

\begin{quote}
\emph{chor ne khazānā zamīn meṅ gāṛā.}
\end{quote}

\begin{quote}
The thief buried the treasure in the ground.
\end{quote}

Past, with \ur{نے} (he or she …): \ur{اس} \ur{نے} \ur{بیج} \ur{گاڑا}, \emph{us ne bīj gāṛā}.

\begin{cousinweb}
\textbf{\ur{گاڑنا}} (\emph{gāṛnā}) is inherited: CDIAL lists the Hindi cousin \emph{gāṛnā}, \textquotedblleft{}to bury, to drive in\textquotedblright{}, under a form it reconstructs (CDIAL 3979).
\end{cousinweb}

\begin{grammarlens}[title={What you've built}]
One more word for telling what happened.
\end{grammarlens}

\subsection*{Your turn}
\begin{itemize}
\raggedright
  \item \emph{Say it:} \emph{gāṛnā}
  \item \emph{Say it:} \emph{gāṛnā}, once more
  \item \emph{Say it:} \emph{chor ne khazānā zamīn meṅ gāṛā.} — The thief buried the treasure in the ground.
  \item \emph{Recall:} say the Urdu for torment, then the Urdu for a treasure, then say \emph{gāṛnā} again
\end{itemize}

\begin{tcolorbox}[breakable,colback=teal!4,colframe=teal!35!black,title={Before you move on}]
What does \ur{گاڑنا} mean? (\textquotedblleft{}To bury\textquotedblright{}.) Say it once more.
\end{tcolorbox}
\section[\texorpdfstring{\ur{ناگ} (nāg)}{ناگ (nāg)}]{\ur{ناگ} (nāg) — a cobra; a serpent}

\label{lesson:UR-C327-nag}



\begin{quote}
Before the new one: say the Urdu for a treasure, then the Urdu for to bury.
\end{quote}

\subsection*{You'll want to know: \ur{ناگ}}
\textbf{\ur{ناگ}} — \emph{nāg} — \textquotedblleft{}a cobra; a serpent\textquotedblright{}.

\begin{quote}
\ur{خزانے} \ur{پر} \ur{ایک} \ur{ناگ} \ur{بیٹھا} \ur{تھا۔}
\end{quote}

\begin{quote}
\emph{khazāne par ek nāg baiṭhā thā.}
\end{quote}

\begin{quote}
A cobra sat on the treasure.
\end{quote}

Masculine.

\begin{cousinweb}
is inherited from Sanskrit \emph{nāga}, \textquotedblleft{}a serpent\textquotedblright{} (Monier-Williams). In tales the \textbf{\ur{ناگ}} guards treasure.
\end{cousinweb}

\begin{grammarlens}[title={What you've built}]
One more word for telling what happened.
\end{grammarlens}

\subsection*{Your turn}
\begin{itemize}
\raggedright
  \item \emph{Say it:} \emph{nāg}
  \item \emph{Say it:} \emph{nāg}, once more
  \item \emph{Say it:} \emph{khazāne par ek nāg baiṭhā thā.} — A cobra sat on the treasure.
  \item \emph{Recall:} say the Urdu for a treasure, then the Urdu for to bury, then say \emph{nāg} again
\end{itemize}

\begin{tcolorbox}[breakable,colback=teal!4,colframe=teal!35!black,title={Before you move on}]
What does \ur{ناگ} mean? (\textquotedblleft{}A cobra; a serpent\textquotedblright{}.) Say it once more.
\end{tcolorbox}
\section[\texorpdfstring{\ur{ڈسنا} (ḍasnā)}{ڈسنا (ḍasnā)}]{\ur{ڈسنا} (ḍasnā) — to bite (of a snake)}

\label{lesson:UR-C327-dasna}



\begin{quote}
Before the new one: say the Urdu for to bury, then the Urdu for a cobra.
\end{quote}

\subsection*{You'll want to know: \ur{ڈسنا}}
\textbf{\ur{ڈسنا}} — \emph{ḍasnā} — \textquotedblleft{}to bite (of a snake)\textquotedblright{}.

\begin{quote}
\ur{ناگ} \ur{نے} \ur{چور} \ur{کو} \ur{ڈسا۔}
\end{quote}

\begin{quote}
\emph{nāg ne chor ko ḍasā.}
\end{quote}

\begin{quote}
The cobra bit the thief.
\end{quote}

Past, with \ur{نے} (he or she …): \ur{اس} \ur{نے} \ur{اسے} \ur{ڈسا}, \emph{us ne use ḍasā}.

\begin{cousinweb}
\textbf{\ur{ڈسنا}} (\emph{ḍasnā}) is inherited: CDIAL lists the Hindi cousin \emph{ḍasnā}, \textquotedblleft{}bites, stings\textquotedblright{}, under Sanskrit \emph{dáśati} (CDIAL 6230).
\end{cousinweb}

\begin{grammarlens}[title={What you've built}]
One more word for telling what happened.
\end{grammarlens}

\subsection*{Your turn}
\begin{itemize}
\raggedright
  \item \emph{Say it:} \emph{ḍasnā}
  \item \emph{Say it:} \emph{ḍasnā}, once more
  \item \emph{Say it:} \emph{nāg ne chor ko ḍasā.} — The cobra bit the thief.
  \item \emph{Recall:} say the Urdu for to bury, then the Urdu for a cobra, then say \emph{ḍasnā} again
\end{itemize}

\begin{tcolorbox}[breakable,colback=teal!4,colframe=teal!35!black,title={Before you move on}]
What does \ur{ڈسنا} mean? (\textquotedblleft{}To bite (of a snake)\textquotedblright{}.) Say it once more.
\end{tcolorbox}
\section[\texorpdfstring{\ur{ہیرا} (hīrā)}{ہیرا (hīrā)}]{\ur{ہیرا} (hīrā) — a diamond}

\label{lesson:UR-C327-hira}



\begin{quote}
Before the new one: say the Urdu for a cobra, then the Urdu for to bite (of a snake).
\end{quote}

\subsection*{You'll want to know: \ur{ہیرا}}
\textbf{\ur{ہیرا}} — \emph{hīrā} — \textquotedblleft{}a diamond\textquotedblright{}.

\begin{quote}
\ur{خزانے} \ur{میں} \ur{سو} \ur{ہیرے} \ur{تھے۔}
\end{quote}

\begin{quote}
\emph{khazāne meṅ sau hīre the.}
\end{quote}

\begin{quote}
There were a hundred diamonds in the treasure.
\end{quote}

Masculine.

\begin{cousinweb}
\textbf{\ur{ہیرا}} (\emph{hīrā}) is inherited from Sanskrit \emph{hīra}, \textquotedblleft{}a diamond\textquotedblright{} (CDIAL 14130), where CDIAL lists the Hindi cousin \emph{hīrā}.
\end{cousinweb}

\begin{grammarlens}[title={What you've built}]
One more word for telling what happened.
\end{grammarlens}

\subsection*{Your turn}
\begin{itemize}
\raggedright
  \item \emph{Say it:} \emph{hīrā}
  \item \emph{Say it:} \emph{hīrā}, once more
  \item \emph{Say it:} \emph{khazāne meṅ sau hīre the.} — There were a hundred diamonds in the treasure.
  \item \emph{Recall:} say the Urdu for a cobra, then the Urdu for to bite (of a snake), then say \emph{hīrā} again
\end{itemize}

\begin{tcolorbox}[breakable,colback=teal!4,colframe=teal!35!black,title={Before you move on}]
What does \ur{ہیرا} mean? (\textquotedblleft{}A diamond\textquotedblright{}.) Say it once more.
\end{tcolorbox}
